% ===================================================================
% FST Programme Overview:
% Functional Stability Theory across Fourteen Levels of Organisation
% Target: Foundations of Physics / Entropy
% Version: 1.0
% ===================================================================

\documentclass[12pt,a4paper]{article}

% Packages
\usepackage[utf8]{inputenc}
\usepackage[T1]{fontenc}
\usepackage{amsmath,amssymb,amsthm}
\usepackage{graphicx}
\usepackage{hyperref}
\usepackage{booktabs}
\usepackage{array}
\usepackage{tabularx}
\usepackage[margin=2.5cm]{geometry}
\usepackage{natbib}
\usepackage{cleveref}
\usepackage{enumitem}

\hypersetup{
  colorlinks  = true,
  linkcolor   = black,
  citecolor   = blue!60!black,
  urlcolor    = blue!70!black,
  pdfauthor   = {Lukas Geiger},
  pdftitle    = {FST Programme Overview},
}

% Theorem environments
\newtheorem{proposition}{Proposition}
\newtheorem{conjecture}{Conjecture}
\newtheorem{definition}{Definition}
\theoremstyle{remark}
\newtheorem{remark}{Remark}

% Title
\title{The FST Programme:\\
Functional Stability Theory across Fourteen Levels\\
of Natural Organisation}

\author{Lukas Geiger\\
\textit{Independent Researcher, Bernau, Germany}\\
ORCID: \url{https://orcid.org/0009-0005-7296-1534}}

\date{March 2026}

\begin{document}

\maketitle

\begin{abstract}
\noindent
Functional Stability Theory (FST) proposes that the structures observed in nature---from quantum fields to civilisations---are Nash equilibria of thermodynamic games at each scale, where the equilibria of each level constrain and enable the games at the next.
This programme overview maps the FST research agenda onto a fourteen-level hierarchy of natural organisation spanning two branches: the \emph{inorganic branch} (Levels~1--5: quantum fields, elementary particles, atoms, molecules, macromolecular systems) and the \emph{biological branch} (Levels~6--14: amino acids, RNA/DNA, genes, proteins, cells, tissues, individuals, groups, societies).
We identify the cross-cutting patterns that recur across the hierarchy---symmetry breaking as emergent structure, the universal cooperation--defection tension, nested constraints, enforcement escalation, and the cooperation ratchet---and show how each is formalised within the existing FST papers (FST-I through FST-III) or constitutes an open programme for future work (FST-IV through FST-VI).
The Renormalized Free-Energy Principle (RFEP) provides the theoretical foundation, with ``hierarchical selection'' manifested as the eleven-step constraint chain from the Higgs vacuum to societal institutions.
No theological or metaphysical content is included; the framework is purely natural-scientific.

\medskip
\noindent\textbf{Keywords:} functional stability theory, Nash equilibrium, fractal hierarchy, self-organisation, symmetry breaking, cooperation, multilevel selection, entropy production, RFEP
\end{abstract}

\newpage
\setcounter{page}{1}

% ===================================================================
\section{Introduction}
\label{sec:introduction}
% ===================================================================

The Functional Stability Theory (FST) research programme applies game-theoretic optimisation principles---grounded in thermodynamic entropy production and Nash equilibrium---to the full hierarchy of natural organisation. The programme comprises:

\begin{itemize}[leftmargin=2em]
  \item \textbf{FST-I} \citep{Geiger2026a}: Thermodynamic Game Theory of Fundamental Particles (Levels~1--2).
  \item \textbf{FST-II} \citep{Geiger2026b}: Chemical Game Theory---prebiotic chemistry as thermodynamic Nash equilibrium (Levels~3--5).
  \item \textbf{FST-III} \citep{Geiger2026c}: Biological Game Theory---protein folding, gene regulation, and cancer as Nash-frustrated systems (Level~9, with extensions to Levels~8 and~10).
  \item \textbf{RFEP} \citep{GeigerRFEP2026}: The Renormalized Free-Energy Principle, providing the theoretical foundation for hierarchical selection across scales.
\end{itemize}

This document serves as a programme overview. It extracts the structural regularities visible only when the full fourteen-level hierarchy is considered as a whole, maps each level to the corresponding FST paper, and identifies the open programme for levels not yet covered quantitatively.

\subsection{Scope and Method}

The fourteen levels are organised into two branches:

\begin{quote}
\textbf{Inorganic branch} (Levels~1--5): Quantum fields $\to$ Elementary particles $\to$ Atoms $\to$ Molecules $\to$ Macromolecular/mineral systems.

\textbf{Biological branch} (Levels~6--14): Amino acids $\to$ RNA/DNA $\to$ Genes $\to$ Proteins $\to$ Cells $\to$ Tissues $\to$ Individuals $\to$ Groups $\to$ Societies.
\end{quote}

At each level, we identify the \emph{players}, the \emph{strategies}, the \emph{payoff structure}, the \emph{Nash equilibria}, and the \emph{enforcement mechanisms}. The mathematical structure is identical across all levels: entities with configuration-dependent outcomes settle into mutually stable states from which no component can unilaterally improve.

The junction between the two branches---where complex chemistry gives rise to self-replicating molecules---is the origin of life: arguably the most profound Nash equilibrium transition in the history of the universe.


% ===================================================================
\section{The Fourteen-Level Hierarchy}
\label{sec:hierarchy}
% ===================================================================

Table~\ref{tab:hierarchy} presents the complete hierarchy with the mapping to FST papers. Levels marked ``Open Programme'' have been characterised qualitatively but not yet formalised with quantitative FST analysis.

\begin{table}[htbp]
  \centering
  \caption{The fourteen levels of natural organisation and their FST mapping.}
  \label{tab:hierarchy}
  \small
  \begin{tabularx}{\textwidth}{c c l X l}
    \toprule
    \textbf{Level} & \textbf{Branch} & \textbf{System} & \textbf{Characteristic Equilibrium} & \textbf{FST Paper} \\
    \midrule
    1  & Inorganic  & Quantum fields       & Vacuum state, Higgs mechanism           & FST-I \\
    2  & Inorganic  & Elementary particles  & Colour confinement, proton stability    & FST-I \\
    3  & Inorganic  & Atoms                & Iron peak, noble gas configurations     & FST-II \\
    4  & Inorganic  & Molecules            & Molecular geometry, chemical equilibrium & FST-II \\
    5  & Inorganic  & Macromolecular sys.  & Crystals, lipid bilayers, autocatalysis & FST-II \\
    \midrule
    6  & Biological & Amino acids          & Genetic code as error-robust assignment & FST-III \\
    7  & Biological & RNA/DNA              & Eigen's error threshold, Watson-Crick   & FST-III \\
    8  & Biological & Genes                & Mendelian segregation, genome cooperation & FST-III \\
    9  & Biological & Proteins             & Native fold as Nash equilibrium         & FST-III \\
    10 & Biological & Cells                & Multicellular cooperation vs.\ cancer   & FST-III \\
    11 & Biological & Tissues              & Homeostasis, host-microbiome mutualism  & FST-III \\
    \midrule
    12 & Biological & Individuals          & Hawk-Dove, ESS, Tit-for-Tat            & \textit{Open: FST-IV?} \\
    13 & Biological & Groups               & Multilevel selection, eusociality       & \textit{Open: FST-V?} \\
    14 & Biological & Societies            & Social contract, institutional design   & \textit{Open: FST-VI?} \\
    \bottomrule
  \end{tabularx}
\end{table}


% ===================================================================
\section{Cross-Cutting Pattern: Symmetry Breaking as Emergent Structure}
\label{sec:symmetry_breaking}
% ===================================================================

At every inorganic level, the system transitions from a high-symmetry (disordered) state to a lower-symmetry (structured) equilibrium. This is not analogous to defection---it is closer to \emph{cooperation}: the components collectively settle on a shared lower-symmetry configuration, sacrificing individual freedom (symmetric equivalence) for collective stability.

\begin{table}[htbp]
  \centering
  \caption{Symmetry breaking as emergent cooperative structure across the inorganic branch.}
  \label{tab:symmetry_breaking}
  \begin{tabularx}{\textwidth}{c l X X}
    \toprule
    \textbf{Level} & \textbf{System} & \textbf{Symmetric State (Disordered)} & \textbf{Broken-Symmetry Equilibrium (Structured)} \\
    \midrule
    1 & Quantum fields & Symmetric vacuum ($\langle\phi\rangle = 0$) & Higgs mechanism ($\langle\phi\rangle = v$) \\
    2 & Particles & Matter-antimatter symmetry & Baryogenesis (matter dominance) \\
    3 & Atoms & Spherical electron cloud & Shell structure, orbital shapes \\
    4 & Molecules & Free atoms (maximal symmetry) & Molecular geometry (reduced symmetry) \\
    5 & Macrosystems & Homogeneous solution & Crystal lattice, membrane, autocatalytic network \\
    \bottomrule
  \end{tabularx}
\end{table}

\paragraph{Connection to FST Pattern~A.}
In the FST framework, the disorder-to-structure transition at each inorganic level corresponds to the \emph{Second-Order Resolvent Dominance} mechanism formalised in FST-I \citep{Geiger2026a}: the broken-symmetry state is the resolvent-stable fixed point where the Hessian (second-order perturbation analysis) yields strictly positive eigenvalues in all non-gauge directions. The symmetric state is a saddle point---an unstable Nash equilibrium---while the broken-symmetry state is the stable Nash equilibrium.


% ===================================================================
\section{Cross-Cutting Pattern: The Universal Prisoner's Dilemma}
\label{sec:prisoners_dilemma}
% ===================================================================

At every biological level, entities face a variant of the Prisoner's Dilemma: cooperation yields collective benefit, but individual defection yields short-term advantage at the expense of the group.

\begin{table}[htbp]
  \centering
  \caption{The cooperation--defection tension at each level of biological organisation.}
  \label{tab:pd_levels}
  \begin{tabularx}{\textwidth}{c l X X}
    \toprule
    \textbf{Level} & \textbf{Players} & \textbf{Cooperation} & \textbf{Defection} \\
    \midrule
    6  & Amino acids  & Stable chemical interactions         & Destabilising mismatches \\
    7  & RNA/DNA      & Faithful replication                 & Parasitic replication \\
    8  & Genes        & Fair segregation (Mendelian)          & Meiotic drive \\
    9  & Proteins     & Native fold                          & Amyloid aggregation \\
    10 & Cells        & Differentiation, growth control      & Cancer \\
    11 & Tissues      & Homeostatic contribution             & Hypertrophy, inflammation \\
    12 & Individuals  & Reciprocal altruism                  & Cheating, aggression \\
    13 & Groups       & Collective action                    & Free-riding \\
    14 & Societies    & Institutional compliance             & Corruption, war \\
    \bottomrule
  \end{tabularx}
\end{table}

\paragraph{Connection to FST.}
The cooperation--defection tension is quantified in FST-II through replicator dynamics \citep{Geiger2026b} and in FST-III through Nash frustration in protein folding landscapes \citep{Geiger2026c}. In FST-III, the native fold of a protein is shown to be a Nash equilibrium from which no single residue can deviate to lower free energy, while amyloid aggregation represents a pathological counter-equilibrium---a ``bad'' Nash equilibrium stabilised by $\beta$-sheet-rich cross-interactions.


% ===================================================================
\section{Cross-Cutting Pattern: Nested Constraints}
\label{sec:nested_constraints}
% ===================================================================

The Nash equilibria of lower levels are not merely analogous to those of higher levels---they are \emph{constitutive}. Each level's equilibria constrain and enable the games at the next level. This eleven-step constraint chain runs as follows:

\begin{enumerate}[label=(\arabic*)]
  \item The Higgs vacuum (Level~1) determines particle masses (Level~2).
  \item Particle properties (Level~2) determine nuclear and atomic structure (Level~3).
  \item The periodic table (Level~3) determines molecular bonding possibilities (Level~4).
  \item Molecular properties (Level~4) determine macromolecular self-assembly (Level~5).
  \item Macromolecular systems (Level~5) set the boundary conditions for biological evolution (Level~6).
  \item The genetic code (Level~6) constrains gene evolution (Level~8).
  \item The protein folding landscape (Level~9) constrains cellular behaviour (Level~10).
  \item Cellular cooperation (Level~10) enables tissue-level organisation (Level~11).
  \item Tissue homeostasis (Level~11) supports individual organisms (Level~12).
  \item Individual strategies (Level~12) shape group dynamics (Level~13).
  \item Group norms (Level~13) form the basis of societal institutions (Level~14).
\end{enumerate}

This is the precise sense in which the hierarchy is \emph{fractal}: the structure at each scale is shaped by, and recapitulates, the structure at the scale below.

\paragraph{Connection to RFEP.}
The Renormalized Free-Energy Principle \citep{GeigerRFEP2026} formalises this nesting as ``hierarchical selection'': each level's free-energy landscape is constrained by the equilibria of the level below, and the renormalization group (RG) provides the mathematical machinery for the scale transition. The constraint chain above is the empirical manifestation of the RFEP's hierarchical selection principle.

\paragraph{Quantitative confirmation.}
FST-I provides the first quantitative test: the observed fine-structure constant $\alpha \approx 1/137$ achieves 98.8\% of the maximum integrated entropy production in a semi-analytic stellar model \citep{Geiger2026a}. This result demonstrates that the Level~1--2 equilibrium (Standard Model parameters) is not arbitrary but entropy-optimal---consistent with the nested-constraint logic that physical constants at lower levels are selected to maximise dissipation cascading through all higher levels. The ratio $m_e/m_p \approx 1/1836$ receives a similar entropy-dominance interpretation within the same framework.


% ===================================================================
\section{Cross-Cutting Pattern: Enforcement Escalation}
\label{sec:enforcement}
% ===================================================================

Each level develops enforcement mechanisms of increasing sophistication. A gradient emerges across the full hierarchy, from inviolable physical law at the lowest levels to contingent institutional design at the highest:

\begin{table}[htbp]
  \centering
  \caption{Enforcement mechanisms across all fourteen levels.}
  \label{tab:enforcement}
  \begin{tabularx}{\textwidth}{c X X}
    \toprule
    \textbf{Level} & \textbf{Primary Enforcement} & \textbf{Character} \\
    \midrule
    1  & Conservation laws, gauge symmetry                         & Absolute (no violation possible) \\
    2  & Colour confinement, quantum number conservation           & Absolute \\
    3  & Pauli exclusion, Coulomb barrier                          & Absolute (quantum) + statistical \\
    4  & Activation energy barriers, thermodynamics                & Statistical (can be overcome) \\
    5  & Compartmental boundaries, catalytic specificity           & Functional (can be disrupted) \\
    6  & tRNA synthetase specificity                               & Molecular policing \\
    7  & DNA polymerase proofreading, Watson-Crick thermodynamics  & Molecular + thermodynamic \\
    8  & Meiotic machinery, transposon silencing                   & Molecular policing \\
    9  & Chaperones, ubiquitin-proteasome, autophagy               & Molecular surveillance \\
    10 & Tumour suppressors, immune surveillance, apoptosis        & Cellular surveillance \\
    11 & Hormonal regulation, negative feedback loops              & Systemic regulation \\
    12 & Territoriality, dominance hierarchies, reciprocity        & Behavioural \\
    13 & Altruistic punishment, ostracism, reputation              & Social-behavioural \\
    14 & Legal systems, markets, religion, democracy               & Institutional \\
    \bottomrule
  \end{tabularx}
\end{table}


% ===================================================================
\section{Cross-Cutting Pattern: The Cooperation Ratchet}
\label{sec:cooperation_ratchet}
% ===================================================================

Each successful resolution of a cooperation problem enables the emergence of a new level of organisation, which in turn faces its own cooperation problem. This is the ``major transitions in evolution'' framework of \citet{MaynardSmith1995}, extended to the inorganic domain and recast in explicitly game-theoretic terms:

\begin{quote}
Quantum field equilibria $\to$ stable particles $\to$ atoms $\to$ molecules $\to$ macromolecular systems $\to$ amino acid cooperation $\to$ stable proteins $\to$ cellular machinery $\to$ multicellular organisms $\to$ social groups $\to$ civilisations.
\end{quote}

Each major transition is the resolution of a strategic dilemma at one level, creating the players for the game at the next level. The ratchet is irreversible in the sense that higher-level players cannot exist without the equilibria of lower levels being maintained.

\paragraph{FST interpretation.}
In the FST framework, each step of the cooperation ratchet corresponds to a Kadanoff block-spin step in the renormalization group: the players of the lower level are ``blocked'' into a super-player at the next level. The game logic (cooperation vs.\ defection, stability vs.\ disruption) is preserved across the block step---this is the structural self-similarity that justifies the term ``fractal'' in the programme description.

The quantitative FST papers provide explicit realisations:
\begin{itemize}[leftmargin=2em]
  \item \textbf{FST-I:} The transition from quantum fields (Level~1) to stable particles (Level~2) is formalised as a symmetry-breaking Nash equilibrium with the Higgs vacuum expectation value as the order parameter.
  \item \textbf{FST-II:} The transition from molecules (Level~4) to autocatalytic networks (Level~5) is formalised through replicator dynamics and bistability analysis.
  \item \textbf{FST-III:} The transition within biological organisation (Level~9: protein folding) is formalised through Nash frustration scores and gradient-flow analysis on the free-energy landscape.
\end{itemize}


% ===================================================================
\section{Open Programme: Levels 12--14}
\label{sec:open_programme}
% ===================================================================

The current FST papers (I--III) cover Levels~1--11 with varying degrees of quantitative detail. Levels~12--14 have been characterised qualitatively (see Section~\ref{sec:prisoners_dilemma}) but remain open for formal FST treatment.

\subsection{Level 12: Individuals --- Toward FST-IV}

Level~12 is the classical domain of evolutionary game theory \citep{MaynardSmith1982}: Hawk-Dove games, Evolutionarily Stable Strategies (ESS), and the Iterated Prisoner's Dilemma with strategies like Tit-for-Tat \citep{Axelrod1984}.

A future \textbf{FST-IV} would connect behavioural game theory to the FST framework by:
\begin{itemize}[leftmargin=2em]
  \item Formalising individual behavioural strategies as mixed-strategy Nash equilibria within the RFEP hierarchy.
  \item Connecting frequency-dependent selection to the replicator dynamics already established in FST-II.
  \item Identifying the entropy-production proxy for behavioural fitness (metabolic rate, information processing capacity).
\end{itemize}

\subsection{Level 13: Groups --- Toward FST-V}

Level~13 addresses social dilemmas: the tension between within-group selection (favouring defectors) and between-group selection (favouring cooperative groups) \citep{Wilson2007}. Eusocial insects represent the extreme equilibrium where individual defection is almost entirely suppressed.

A future \textbf{FST-V} would:
\begin{itemize}[leftmargin=2em]
  \item Formalise multilevel selection as a two-scale game with RG-connected Nash equilibria.
  \item Quantify the enforcement cost of cooperation maintenance (punishment, ostracism) using Nowak's five rules for the evolution of cooperation \citep{Nowak2006}.
  \item Connect Hamilton's rule ($rB > C$) \citep{Hamilton1964} to the RFEP's hierarchical free-energy minimisation.
\end{itemize}

\subsection{Level 14: Societies --- Toward FST-VI}

Level~14 encompasses institutional design: legal systems, markets, democratic accountability, and international institutions. The social contract is interpreted as a Nash equilibrium where citizens accept constraints on individual freedom in exchange for collective security.

A future \textbf{FST-VI} would:
\begin{itemize}[leftmargin=2em]
  \item Formalise different institutional arrangements (democracy, autocracy) as distinct Nash equilibria with different total welfare.
  \item Analyse the Stag Hunt structure of international cooperation (climate agreements, trade treaties).
  \item Connect institutional enforcement (the ``cheapest'' being internalised norms---religion, ideology) to the enforcement escalation gradient of Table~\ref{tab:enforcement}.
\end{itemize}


% ===================================================================
\section{Synthesis: Self-Similar Structure Across Scales}
\label{sec:synthesis}
% ===================================================================

The five cross-cutting patterns identified in Sections~\ref{sec:symmetry_breaking}--\ref{sec:cooperation_ratchet} reveal a striking property of the FST programme: \emph{fractal self-similarity}. Pattern~A (Second-Order Resolvent Dominance) recurs at each organisational level with structural invariance---a hallmark of fractal organisation in the mathematical sense (self-similar structure under scale transformation). The term ``fractal'' is used here as a \emph{descriptive adjective} for this observed self-similarity, not as part of the programme name.

\begin{enumerate}
  \item \textbf{Two motifs, one grammar.} The inorganic branch is governed by the tension between symmetry and symmetry-breaking; the biological branch adds the tension between cooperation and defection. Both share the same underlying logic: entities with interdependent outcomes settle into mutually stable configurations.

  \item \textbf{From absolute to contingent enforcement.} The enforcement hierarchy transitions from inviolable physical law (Levels~1--2) through statistical barriers (Levels~3--4) and functional boundaries (Level~5) to molecular policing (Levels~6--9), cellular surveillance (Levels~10--11), behavioural enforcement (Levels~12--13), and institutional enforcement (Level~14).

  \item \textbf{The equilibrium-to-dissipation transition.} The most consequential Nash equilibrium shift occurs at the boundary between the two branches: the transition from thermodynamic equilibrium-seeking systems to far-from-equilibrium self-maintaining systems. This is the origin of life---not as a violation of physics, but as a new class of physical Nash equilibrium.

  \item \textbf{Nested equilibria.} The Nash equilibria of each level form the boundary conditions for the games at the next level, creating the eleven-step constraint chain of Section~\ref{sec:nested_constraints}.

  \item \textbf{The cooperation ratchet.} Each major transition in evolution is the resolution of a strategic dilemma at one level, creating the players for the game at the next.
\end{enumerate}

\paragraph{Quantitative status.}
FST-I demonstrates entropy dominance of Standard Model parameters (0.47 normalised entropy production for $\alpha$, where the observed value achieves 98.8\% of the theoretical maximum). FST-II demonstrates replicator-dynamic bistability in prebiotic autocatalytic systems. FST-III demonstrates Nash frustration as a quantitative predictor of protein misfolding pathology. Together, these results provide quantitative confirmation of the FST programme across Levels~1--11. Levels~12--14 remain qualitatively characterised and constitute the open programme.


% ===================================================================
\section{Conclusion}
\label{sec:conclusion}
% ===================================================================

Nature organises itself through nested games. From the quantum vacuum to civilisation, the same structural pattern recurs: entities with configuration-dependent outcomes settling into mutually stable states---Nash equilibria---from which no component can unilaterally improve.

The FST programme provides a unified language for this hierarchy. The existing papers (FST-I through FST-III) cover Levels~1--11 with quantitative analysis; the RFEP provides the theoretical foundation via hierarchical selection and renormalization-group scale transitions. Levels~12--14 (individuals, groups, societies) constitute the open programme, where behavioural game theory, multilevel selection theory, and institutional design await formal integration into the FST framework.

The single lesson of the programme is this: the structures we observe---particles, atoms, molecules, cells, organisms, societies---are the Nash equilibria of thermodynamic games at each scale, nested within each other from the Planck scale to geopolitics. The cooperation ratchet that has driven the hierarchy upward---repeatedly, at every scale, against the relentless pressure of defection and disorder---is the central phenomenon that the FST programme seeks to formalise and quantify.


% ===================================================================
% AI Disclosure
% ===================================================================
\section*{AI Disclosure (Level 5)}

This document was developed with extensive intellectual assistance from Claude (Anthropic). The AI contributed to structural organisation, literature synthesis, and text generation. The author provided the research direction, selected the content, and reviewed all material for accuracy.


% ===================================================================
% References
% ===================================================================
\newpage
\begin{thebibliography}{99}

\bibitem[Axelrod(1984)]{Axelrod1984}
  Axelrod, R. (1984).
  \textit{The Evolution of Cooperation}.
  Basic Books, New York.

\bibitem[Geiger(2026a)]{Geiger2026a}
  Geiger, L. (2026).
  Thermodynamic Game Theory of Fundamental Particles: The Standard Model as Nash-Optimal Equilibrium.
  Preprint (FST-I).

\bibitem[Geiger(2026b)]{Geiger2026b}
  Geiger, L. (2026).
  Chemical Game Theory: Prebiotic Chemistry as Thermodynamic Nash Equilibrium under the Maximum Entropy Production Principle.
  Preprint (FST-II).

\bibitem[Geiger(2026c)]{Geiger2026c}
  Geiger, L. (2026).
  Biological Game Theory: Protein Folding, Gene Regulation, and Cancer as Nash-Frustrated Systems.
  Preprint (FST-III).

\bibitem[Geiger(2026d)]{GeigerRFEP2026}
  Geiger, L. (2026).
  The Renormalized Free-Energy Principle: Hierarchical Selection across Scales.
  Zenodo. DOI: 10.5281/zenodo.19039190.

\bibitem[Hamilton(1964)]{Hamilton1964}
  Hamilton, W.D. (1964).
  The genetical evolution of social behaviour.
  \textit{Journal of Theoretical Biology}, 7(1), 1--52.

\bibitem[Maynard~Smith(1982)]{MaynardSmith1982}
  Maynard~Smith, J. (1982).
  \textit{Evolution and the Theory of Games}.
  Cambridge University Press.

\bibitem[Maynard~Smith and Szathm{\'a}ry(1995)]{MaynardSmith1995}
  Maynard~Smith, J. and Szathm{\'a}ry, E. (1995).
  \textit{The Major Transitions in Evolution}.
  Oxford University Press.

\bibitem[Nowak(2006)]{Nowak2006}
  Nowak, M.A. (2006).
  Five rules for the evolution of cooperation.
  \textit{Science}, 314(5805), 1560--1563.

\bibitem[Wilson and Wilson(2007)]{Wilson2007}
  Wilson, D.S. and Wilson, E.O. (2007).
  Rethinking the theoretical foundation of sociobiology.
  \textit{The Quarterly Review of Biology}, 82(4), 327--348.

\end{thebibliography}

% ===================================================================
\end{document}
% ===================================================================
